Metagenomic next-generation sequencing can interrogate complex or unexpected microbial mixtures without a fixed target panel, and it has become an important route for infectious-disease investigation when targeted assays are insufficient \citep{Wilson2014,Wilson2019,Chiu2019}. Yet the reads entering downstream analysis are often short, trimmed, ambiguous and locally corrupted. Adapter removal, quality trimming and low-biomass contamination can all change the effective sequence evidence available to a pipeline before classification begins \citep{Martin2011,Bolger2014,Salter2014}. Read length can affect homolog recovery and taxonomic assignment, although the magnitude and direction of the effect depend on the classifier, database, taxonomic group and error profile \citep{Wommack2008,Pearman2020,Wood2014,Menzel2016}. Reduced context can narrow the redundancy of local sequence evidence and make it harder to determine which information categories remain available. A read representation may already have compressed local composition, biochemical regularity or positional information before a downstream benchmark reports success or failure. Short-read representation therefore defines which categories of evidence remain available to a later classifier, alignment routine or database query.

Many metagenomic workflows rely on exact or near-exact match evidence. Canonical k-mer, minimizer, alignment and database-backed systems such as \textit{Kraken}-family tools, \textit{Centrifuge}, \textit{CLARK}, \textit{Kaiju} and related indexing frameworks have established the practical value of explicit sequence matching for taxonomic assignment and downstream biological interpretation \citep{Wood2014,Wood2019,Ounit2015,Kim2016,Menzel2016}. Community benchmarks further show that performance depends on database coverage, sample composition, novelty and evaluation context, not only on the classifier name \citep{Sczyrba2017,Meyer2022}. Database-linked match evidence remains the appropriate backbone for species discrimination, allele-sensitive interpretation and resistance-gene calling. Spaced-seed indexing extends this line by tuning match patterns for sensitivity under mutation \citep{Brinda2015}. Alignment-free comparison adds a related interpretive problem: similarity scores are affected by sequence length and require objective functions or score-distribution analyses before they can be assigned biological meaning \citep{Swain2022AFScore}. Additional audit coordinates may therefore be especially informative toward the lower end of the 50--75 bp range, where local matching evidence has less redundancy than at 150 bp. These systems identify database matches, whereas the present study characterizes how composition, nucleotide-property-coded and positional signals change within compact representations when reads are short, noisy or partially degraded.

Vectorized k-mer representations provide the first broad alternative to explicit matching. Here the central operation is to collapse each read or fragment into a composition profile, such as count, frequency, relative-frequency, TF-IDF, or background-adjusted features. PlasFlow and the BMC Bioinformatics 2018 bacteria classifier show that this compact route can be highly competitive when the task is composition-driven rather than order-driven \citep{PlasFlow2018,Fiannaca2018}. Resource-saving k-mer classification extends the same logic by showing that low-dimensional k-mer distributions plus classical machine learning can remain competitive \citep{Fuhl2023}. Recent cross-genome analysis of k-mer rarity further shows that k-mer length changes the balance between abundance, compositional bias and taxonomic specificity \citep{Chantzi2024Rarity}. A complementary study used both bulk and positional k-mers in an interpretable mixture model, demonstrating that composition and positional organization can support different sequence analyses \citep{Wang2024LDA}. These studies motivate the present distinction between a compact composition backbone and a separate coarse positional summary. They do not, however, establish how nucleotide-property-coded positional channels behave under matched short-read perturbations. Accordingly, vectorized k-mer methods are treated here as composition-focused comparators whose performance depends on how much local composition survives compression in the short-read regime.

Hashing and sketching form a separate compression family rather than a lower-rank version of singular-value decomposition (SVD). MinHash, locality-sensitive hashing (LSH), histosketch and related strategies map k-mer sets or profiles into fixed-length signatures for rapid similarity search, large-scale database indexing or ultra-lightweight compression \citep{Ondov2016,Zielezinski2017,Kssd2021,HULK2019}. Kssd and the HULK/histosketch line compress representations in service of indexing or neighbourhood search rather than preserving linear variance structure. The present analysis instead evaluates signal retention in compact vectors under fixed probes and perturbations. Sketching is therefore treated as a neighbouring compression paradigm and assessed in its native similarity geometry.

Biochemical and signal-based encodings add a second alternative representation family. Chaos-game and genomic-signal formulations, together with numerical mappings such as Voss, tetrahedron and electron-ion interaction potential (EIIP), expose sequence structure through algebraic or physicochemical coordinates rather than literal token identity alone \citep{Jeffrey1990,Voss1992,Anastassiou2001,Cristea2002,MendizabalRuiz2017,MendizabalRuiz2018GSP,Nair2006}. Nucleotide chemical-property encodings have also represented ring structure, chemical functionality and hydrogen-bond class, often together with cumulative nucleotide density or other position-dependent frequencies \citep{Chen2017iDNA4mC}. These mappings have been used well beyond metagenomic classification, including promoter, regulatory-site, nucleosome-positioning and nucleotide-modification prediction. Their value is therefore not specific to mNGS or to k-mer classifiers. At the same time, the biological interpretation of a named property does not by itself establish that its numerical coordinate is causal, non-redundant or robust across sequence lengths and tasks.

Hybrid handcrafted descriptors already combine composition, physicochemical properties and sequence-order information. Pseudo k-tuple nucleotide composition (PseKNC) augments k-tuple frequencies with lagged correlations between oligonucleotide properties, and PseKNC-General and repDNA expose broad families of autocorrelation, pseudo-composition and user-defined physicochemical descriptors \citep{Chen2014PseKNC,Chen2015PseKNCGeneral,Liu2015repDNA}. Pseudo electron-ion interaction potential (PseEIIP) similarly weights trinucleotide composition by EIIP, whereas platforms such as iFeatureOmega catalogue composition, position-specific, autocorrelation and physicochemical encoders within one feature-engineering system \citep{He2018PseEIIP,Chen2022iFeatureOmega}. Recent multiscale density encodings further show that local and broader position-dependent nucleotide patterns can be combined with learned models \citep{Zabihi2026EDEN}. Learning-based multimodal models provide a related but distinct precedent: MIXBend aligns k-mer sequence features with nucleotide physicochemical and positional information in a self-supervised framework for 50-bp DNA bendability prediction \citep{Yang2024MIXBend}. These studies establish that composition, physicochemical and positional signals can be fused for task-specific prediction. The methodological gap is narrower: such representations are usually assessed through the performance of a task-specific predictor, whereas a common matched-perturbation audit that attributes stability and failure modes to individual blocks remains less developed. Here the contribution is a compact, training-free and block-decomposable representation-and-audit protocol for controlled short-read perturbations, rather than a claim that the underlying nucleotide properties or their fusion are new.

Sequence-token deep learning classifiers make the neighbouring problem space especially relevant. DeepMicrobes introduced deep learning for metagenomic taxonomic classification, using canonical k-mer tokenization, embedding, BiLSTM and attention for read-level assignment \citep{Liang2020}. MetaTransformer extended this family with self-attention and alternative tokenization or embedding schemes, including character-level, BPE, k-mer and hash/LSH variants \citep{Wichmann2023}. DeepMicroClass provides a related example of a task-specific deep classifier evaluated across sequence classes and external CAMI II data, although its primary unit is a metagenomic contig rather than a 50--150 bp read \citep{Hou2024DeepMicroClass}. BERTax showed that a small-k Transformer can support hierarchical taxonomy when pretraining and output heads are aligned to the task \citep{BERTax2022}. BarcodeBERT further demonstrated that barcode-scale biodiversity tasks benefit from domain-specific pretraining and masking choices, although its setting is barcode classification rather than shotgun short-read mNGS \citep{BarcodeBERT2026}. Large pre-trained DNA language models extend this representation family as general-purpose encoders \citep{Zhou2024,DallaTorre2025,Nguyen2023,Nguyen2024Evo,Fishman2025}. Such learned embeddings can provide strong predictive features, particularly when pretraining data and compute are available. These models provide a predictive-performance reference point, whereas the present study examines a low-dimensional, training-free and block-decomposable representation under controlled short-read perturbation.

Raw-sequence and position-aware models complete the picture. DeepVirFinder, Virtifier, and DETIRE demonstrate that CNN or hybrid architectures can extract useful motif- and composition-level signal from viral fragments, but their target is viral detection rather than multi-class short-read taxonomy \citep{DeepVirFinder2020,Virtifier2022,DETIRE2023}. KmerAperture adds an important boundary case by showing that ordinary k-mer set comparisons can lose synteny, whereas matching k-mers back to their original lists can retain positional information about core and accessory differences \citep{KmerAperture2024}. Collectively, the reviewed literature has established composition vectors, physicochemical descriptors, positional k-mers, task-specific hybrid encoders and learned sequence embeddings. These approaches answer different predictive or retrieval questions. A common protocol for measuring how compact composition, nucleotide-property-coded and coarse positional blocks retain, lose or redistribute signal under matched perturbations remains underdeveloped. The present work is therefore positioned as a representation-and-audit contribution that makes those trade-offs explicit, rather than as a replacement for the strongest task-specific predictor in any one application.

Against this background, we treat short-read representation as an audit problem with several distinct endpoints. The unmet methodological need is a fixed-budget protocol that attributes perturbation responses to local composition, nucleotide-property summaries and coarse relative position. We ask how local k-mer-composition retention, nucleotide-property-coded stability, positional readout accessibility and compactness trade off across controlled short-read settings. CK4P-MSP combines reverse-complement canonical 4-mer composition with global nucleotide-property-coded summaries and multi-scale property pooling to produce a block-decomposable audit coordinate system. The resulting working representation is a compact, training-free coordinate system for controlled short-read perturbation audits. Its methodological contribution is a block-level attribution protocol that measures the channels separately under matched perturbations, extending rather than originating composition--property fusion. The protocol is not intended to establish a universal ranking of representations: CK4 remains a reasonable choice when composition-based retrieval or coarse readout is the only objective, whereas CK4P-MSP is useful when stability and structured local-change responses must be inspected by named feature blocks. Full-position matrices provide high-dimensional positional-readout references, whereas the canonical spaced-property (CSP) control tests the transfer of matching-oriented spaced-seed priors to dense read-level representations. We combine complete K/P/MSP ablation with matched perturbation, comparator and external-probe audits. A shared-template sweep samples every integer length from 50 to 75 bp, and the 100, 125 and 150 bp anchors assess whether lower-range trends persist as sequence context increases. The analysis estimates metric-specific conditional contributions and identifies the operating range of each block under a compact, block-attributable feature budget.
